\chapter{Diskussion und Fazit}
\label{kap:diskussion}

\section{Was ist belegt, was ist nur begründet?}
\label{sec:tragend}

Kapitel~\ref{kap:pipeline} hat gezeigt, dass die Pipeline die Anforderungen
A1--A6 erfüllt. Das ist eine Aussage über den Entwurf, nicht über die Wirkung.
Für eine ehrliche Bilanz muss unterschieden werden zwischen Entscheidungen, deren
Nutzen \emph{gemessen} wurde, und Entscheidungen, die lediglich \emph{begründet}
sind:

\emph{Gemessen durch eigene Trainingsläufe} sind der Okklusionsanteil
(Abschnitt~\ref{sec:ablation}) und der P2-Detektionskopf
(Abschnitt~\ref{sec:p2}) -- in beiden Fällen liegt der Effekt jedoch weitgehend
innerhalb des Konfidenzintervalls, belastbar ist die Konsistenz, nicht die
Einzeldifferenz. \emph{Gemessen ohne Neutraining} sind Inferenzauflösung,
NMS-Schwelle, \code{max\_det} und Test-Time-Augmentation; nur der
Auflösungseffekt übersteigt das Konfidenzintervall deutlich. \emph{Kalibriert an
der erzeugten Verteilung} sind Objektgrößenbereich, Ovalisierungsanteil,
Gitterabstände, eine Basisgröße pro Cluster, Tiefenmodell 75/25 und Randregel --
hier trifft die erzeugte Statistik die reale (Tabelle~\ref{tab:kalibrierung}),
ohne dass der Nutzen für die Detektionsgüte separat gemessen wurde.
\emph{Indirekt gestützt} ist das Bündel der Anti-Abkürzungs-Maßnahmen S1--S6: Der
Sim-to-Real-Abstand von 9,6\PP{} ist -- statt eines Vielfachen, wie bei einer
durchlässigen Signatur zu erwarten -- ein Indiz, aber kein Nachweis für einzelne
Maßnahmen. \emph{Begründet, aber nicht belegt} bleiben Matting-Schwelle, Farb-,
Unschärfe- und Rauschangleich, JPEG-Requantisierung, Schlagschatten, Gamma, HSV,
Mulde, globale Bildstufe, Negativbildanteil sowie die Dellen- und
Perspektivanteile.

Von den 36 dokumentierten Transformationen ist damit der Trainingsnutzen von
dreien durch eine Ablation gemessen, sechs weitere sind an einer Realmessung
kalibriert, und die übrigen 27 sind begründet, aber nicht einzeln belegt.

\absatz{Die unbequeme Bilanz.} Der größte gemessene Einzeleffekt ist \emph{keine}
Augmentierungsentscheidung, sondern die Inferenzauflösung: 47\PP{} \mAPfifty{}
zwischen \code{imgsz}~640 und 1600. Alle Augmentierungsentscheidungen zusammen
bewegen sich in der Größenordnung weniger Prozentpunkte. Das relativiert den
Aufwand der Pipeline nicht -- ohne sie gäbe es überhaupt keinen
Trainingsdatensatz --, verschiebt aber die Priorität: Wer eine solche Anwendung
baut, sollte die Auflösungsfrage \emph{zuerst} klären und erst danach die
Feinheiten der Synthese optimieren.

\absatz{Was die beiden Experimente nicht beantworten.} Die Okklusionsablation
zeigt, \emph{dass} mehr Verdeckung hilft, aber nicht, ob deren
\emph{physikalische Plausibilität} nötig ist -- ein Arm mit gleich starker, aber
generischer Verdeckung (etwa Cutout) hätte das entschieden und fehlt. Damit
bleibt die aufwendigste Begründung der Pipeline, echte Rohrschäfte statt gemalter
Balken, unbelegt. Und der P2-Vergleich beruht auf je einem Trainingslauf pro
Variante, sodass die Streuung über Seeds unbekannt bleibt.

\section{Grenzen und nächste Schritte}
\label{sec:grenzen}

\absatz{Grenzen der Pipeline.} \emph{(1)} Nur neun Quellobjekte -- jedes
erscheint rund 15\,100-mal als gelabeltes Objekt. Die photometrische und
geometrische Streuung mildert das, hebt es aber nicht auf; Augmentierung erzeugt
keine neue Information über die Objektklasse. Das ist der Hauptgrund für die
Wahl eines kleinen Modells und der plausibelste Grund dafür, dass 6000 Bilder
gegenüber 3000 keinen Vorteil brachten. \emph{(2)} Nur diffuses Licht im
Quellmaterial: Harte Schlagschatten, Gegenlicht, Nässe und Schnee fehlen
vollständig und werden nur approximativ über die globale photometrische Stufe
angenähert -- die größte verbleibende Domänenlücke, von der Fehleranalyse
bestätigt. \emph{(3)} Reflect-Padding erzeugt Spiegelsymmetrie, wenn ein
Hintergrund für den gewünschten Zoom zu klein ist; eine spiegelsymmetrische
Bildhälfte ist selbst ein erkennbares Merkmal, die Maßnahme gegen Signatur~S1
erzeugt damit eine schwächere neue Auffälligkeit. \emph{(4)} Sieben Werte sind
Annahmen (Dellen, Perspektive, Schlagschatten, vertikale Spiegelung, globale
Farbtemperatur, Bewegungsunschärfe und die Clustergröße des Layoutmodells). \emph{(5)} Die synthetischen Cluster sind
dichter als real -- bewusst als Robustheitsreserve belassen; die umgekehrte
Abweichung wäre gefährlich gewesen. \emph{(6)} Die Schaftrichtung ist an der
gemessenen Säulenrichtung des Clusters kalibriert, nicht an der Rohrrichtung des
Hintergrunds: Sie entsteht als Senkrechte zur Säulenrichtung mit einer Streuung
von $55^\circ$ bis $125^\circ$, während die vier Clean Plates ihre Rohre
zwischen $-1{,}6^\circ$ und $+6{,}6^\circ$ zur Waagerechten haben
(Struktur-Tensor, \code{src/occluders.py}). In 19,1\,\% der Cluster kreuzt der
Schaft die Hintergrundrohre deshalb unter mehr als $60^\circ$; betroffen sind
21,6\,\% aller erzeugten Bilder. Weil alle Rohre eines Stapels parallel liegen,
ist das dieselbe Unstimmigkeit, gegen die die enge Streuung von $\pm11^\circ$
\emph{innerhalb} eines Clusters gesetzt ist -- zwischen Cluster und Hintergrund
wird sie nicht erzwungen. Auch hier geht die Abweichung in Richtung mehr
Variation, ihr Nutzen ist aber nicht geprüft; ein Angleich wäre billig, weil die
dominante Hintergrundrichtung gemessen vorliegt.

\absatz{Grenzen der Bewertung.} Die schwerwiegendste methodische Schwäche liegt
nicht in der Pipeline, sondern in der Bewertung: \emph{je ein Bild} pro Split.
Daraus folgen breite Konfidenzintervalle ($\pm 0{,}04$ bis $\pm 0{,}05$
\mAPfifty{}), sodass Unterschiede unterhalb dieser Größe nicht entscheidbar sind;
beide Bilder zeigen dieselbe Objektart bei ähnlichem Wetter, die Ergebnisse sagen
also nichts über andere Zaunhersteller, Rohrdurchmesser oder Lichtsituationen;
und die Ground Truth selbst ist unvollständig, weshalb die Precision-Werte eine
Untergrenze sind. Hinzu kommt der in Abschnitt~\ref{sec:boxversatz} gemessene
Größenversatz zwischen Generatorlabel und Handannotation von 1,2 bis 2,5\px:
Er macht \mAPseven{} und \mAPall{} auf \code{test} zu einem erheblichen Teil zu
einem Maß für die Labelkonvention statt für die Lokalisierungsgüte -- eine
Sensitivitätsrechnung hebt \mAPseven{} von 0,320 auf 0,714. Hinzu kommt, dass \code{val} doppelt verwendet wird -- für die
Modellauswahl \emph{und} für die Wahl der Inferenzauflösung. \code{test} ist
damit ein sauberer Endtest, \code{val} aber kein unabhängiger Bewertungssplit
mehr; die belastbaren Zahlen sind die \code{test}-Zahlen.

\absatz{Nächste Schritte.} Jede der genannten Grenzen hat eine Gegenmaßnahme;
nach erwartetem Nutzen geordnet: \emph{(0)} \textbf{Labelkonvention angleichen.}
Eine geänderte Boxschwelle im Generator genügt, und die Sensitivitätsrechnung
lässt auf \code{test} rund 39 Punkte \mAPseven{} erwarten -- der billigste große
Gewinn der Liste, zu erledigen vor jeder weiteren Optimierung.
\emph{(1)} \textbf{Mehr Bewertungsdaten.} Zehn
bis zwanzig annotierte Realbilder würden die Konfidenzintervalle auf ein Maß
bringen, bei dem beide Experimente entscheidbar wären -- mit Abstand der
wirksamste nächste Schritt. \emph{(2)}
\textbf{Leave-one-out-Ablation der Anti-Abkürzungs-Maßnahmen.} Je ein Lauf mit
abgeschaltetem Farb-, Rausch- und Unschärfeangleich sowie ohne JPEG-Durchlauf
würde erstmals belegen, welche Signatur tatsächlich ausgenutzt würde; Zielgröße
wäre der Sim-to-Real-Abstand, denn eine durchlässige Signatur müsste \code{sval}
nach oben und die realen Splits nach unten bewegen. \emph{(3)} \textbf{Ergänzende
Kontrollarme} -- eine labelzahl-kontrollierte Okklusionsablation mit einem
zusätzlichen Arm für generische statt physikalischer Verdeckung, und der direkte
Vergleich von P2-Kopf und kachelweiser Inferenz \parencite{akyon2022sahi} bei
gleichem Speicherbudget. \emph{(4)} \textbf{Breiteres Quellmaterial} -- zusätzliche
Clean Plates bei harter Sonne, Gegenlicht und Nässe gegen die Lücke aus
Punkt~\emph{(2)} der Pipelinegrenzen; die über 150 weiteren Objektausschnitte,
die aus den vorhandenen Realbildern zu gewinnen wären, stehen dagegen im
Zielkonflikt mit Punkt~\emph{(1)} dieser Liste.

\section{Fazit}

Ein ausschließlich auf synthetischen Bildern trainierter Detektor zählt auf zwei
echten, nie im Training gesehenen Aufnahmen mit einem \textbf{Fehler von
2,3\,\%} -- 86 statt 88 und 84 statt 86 Objekte. Für eine Anwendung, für die kein
einziges handannotiertes Trainingsbild existiert, ist das ein brauchbares
Ergebnis. Die zugehörigen Detektionsgüten sind \mAPfifty{} 0,886 und 0,888, bei
geeigneter Inferenzauflösung 0,928 und mit einem feineren Detektionskopf 0,947 --
letzterer allerdings bei \emph{schlechterem} Zählfehler
(Abschnitt~\ref{sec:p2zaehlung}).

Methodisch tragen fünf Einsichten:

\begin{enumerate}[itemsep=2pt]
  \item \textbf{Bei rein synthetischen Trainingsdaten ist die Pipeline die
  Hypothese.} Jede Entscheidung über eine Transformation ist eine Aussage
  darüber, welche Invarianzen in der Zieldomäne gelten. Die Begründungspflicht --
  Messung, physikalische Randbedingung oder ausgewiesene Annahme -- ist deshalb
  keine Formalität, sondern das eigentliche Handwerk.
  \item \textbf{Die dominante Fehlerquelle sind systematische Versätze, nicht
  fehlende Variation.} Farb-, Schärfe-, Rausch- und Kompressionsversatz zwischen
  eingefügten und originalen Pixeln lassen sich nicht durch mehr Zufall beheben;
  sie müssen einzeln geschlossen werden. Der gemessene Sim-to-Real-Abstand von
  9,6\PP{} -- klein genug zum Funktionieren, groß genug als Beleg gegen reines
  Artefaktlernen -- ist das beste verfügbare Indiz, dass dies gelungen ist.
  \item \textbf{Verdeckung ist bei dicht gepackten Kleinobjekten kein Beiwerk.}
  Stärkere physikalisch plausible Verdeckung verbessert das Ergebnis monoton --
  und zwar gegen einen Verlust von einem Drittel der Labels.
  \item \textbf{Auflösung ist die dominierende Größe, auf drei Ebenen:}
  Objektgrößenbereich des Generators, Inferenzauflösung (Faktor zwei in der
  Detektionsgüte) und Wahl des Detektionskopfes. Die Rechnung
  $\mathbb{E}[N_{\text{Zellen}}] = \pi d^2/(4S^2)$ verbindet alle drei -- sie sagt
  vorher, dass ein zusätzlicher Kopf auf Schrittweite~4 genau dort wirkt, wo die
  Objekte kleiner als etwa zwei Gitterzellen sind, und die Messung bestätigt das
  mit 4,6 gegenüber 0,4\PP{} Vorsprung.
  \item \textbf{Die Labelkonvention ist Teil des Messinstruments.} Ein
  systematischer Größenversatz von zwei Pixeln zwischen Generatorlabel und
  Handannotation kostet auf dem Testbild 39 Punkte \mAPseven{} -- mehr als jede
  Augmentierungsentscheidung einbringt. Wer Lokalisierungsmetriken bei
  20\px-Objekten berichtet, muss die Konvention beider Seiten zuerst
  gegeneinander prüfen.
\end{enumerate}
